\chapter{Aldosterone}

\section*{What Is This Test?}
Aldosterone is the primary mineralocorticoid hormone produced by the zona glomerulosa of the adrenal cortex. It is synthesized from cholesterol through a series of enzymatic reactions culminating with aldosterone synthase (CYP11B2), which converts corticosterone to aldosterone. Aldosterone secretion is primarily regulated by the renin-angiotensin-aldosterone system (RAAS) and serum potassium concentration, with ACTH playing a minor regulatory role. Aldosterone's principal physiological action is sodium retention and potassium and hydrogen ion excretion in the distal convoluted tubule and collecting duct of the kidney. Aldosterone binds to the mineralocorticoid receptor (MR) in the cytoplasm, leading to increased synthesis of epithelial sodium channels (ENaC) and Na+/K+-ATPase, promoting Na+ reabsorption and K+ and H+ secretion. The net effect is expansion of extracellular fluid volume, increased blood pressure, and maintenance of potassium homeostasis. Clinically, the aldosterone test (in conjunction with plasma renin) is used to screen for primary aldosteronism (Conn syndrome), the most common identifiable and potentially curable cause of secondary hypertension, with an estimated prevalence of 5--10\% in unselected hypertensive patients and up to 20\% in resistant hypertension \cite{monticone2017prevalence}.

\section*{Alternative Names}
Serum Aldosterone; Plasma Aldosterone; Aldosterone, Serum; ALD; 24-Hour Urinary Aldosterone; Aldosterone-to-Renin Ratio (ARR); Saline Suppression Test; Captopril Challenge Test; Fludrocortisone Suppression Test

\section*{Principle}
Aldosterone is measured using competitive immunoassays on automated analyzers. An anti-aldosterone antibody competes for aldosterone with labeled aldosterone (with a chemiluminescent tag). The labeled aldosterone is bound to the antibody and the chemiluminescent signal is inversely proportional to the aldosterone concentration. Most modern assays use extraction-free (direct) chemiluminescent immunoassays. Common platforms: Roche Cobas ECLIA, Abbott Architect CMIA, Siemens Atellica CLIA, and Diasorin Liaison. However, aldosterone immunoassays are known to have significant inter-assay variability (CV 10--25\%) and cross-reactivity with other steroids (cortisol, corticosterone, and their metabolites). The gold standard reference method is liquid chromatography-tandem mass spectrometry (LC-MS/MS), which offers superior specificity by separating aldosterone from structurally similar steroids, but it is not yet widely available for routine clinical use. Plasma renin is measured for the aldosterone-to-renin ratio (ARR). Renin is measured either as plasma renin activity (PRA, measuring angiotensin I generated per hour at 37 degrees C, reported in ng/mL/hr) by radioimmunoassay, or as direct renin concentration (DRC, measuring immunoreactive renin mass, reported in mU/L or pg/mL) by automated immunometric assay (Roche, Diasorin). The ARR is calculated as aldosterone (ng/dL) divided by renin activity (ng/mL/hr) or by direct renin concentration (mU/L). An ARR >20--30 (using PRA) or >2--3.7 (using DRC) is considered a positive screen for primary aldosteronism, with the exact cutoff varying by assay and laboratory.

\section*{History}
Aldosterone was independently discovered in 1953 by Sylvia Tait and James Tait at the Middlesex Hospital, London, and by Hudson and Wettstein at Ciba in Basel. Its structure was elucidated that same year. The syndrome of primary aldosteronism was described by Jerome W. Conn in 1955 at the University of Michigan in a 34-year-old woman with hypertension, hypokalemia, and an aldosterone-producing adrenal adenoma (now called Conn syndrome) \cite{conn1955primary}. The first radioimmunoassay for plasma aldosterone was developed by Mayes and colleagues in 1970 \cite{mayes1970radioimmunoassay}. The aldosterone-to-renin ratio as a screening test for primary aldosteronism was first proposed by Hiramatsu and colleagues in 1981 \cite{hiramatsu1981screening}. The 2008 Endocrine Society guideline by Funder and colleagues established ARR screening criteria for primary aldosteronism. Subsequent guidelines in 2016 (Funder et al.) refined diagnostic criteria with confirmatory testing using saline suppression, captopril challenge, fludrocortisone suppression, or oral sodium loading \cite{funder2016management}. The recognition that primary aldosteronism is far more common than previously appreciated (prevalence increasing from 1\% to 5--10\% of hypertension as screening expanded) represents one of the major paradigm shifts in hypertension management over the last two decades \cite{young2007primary, mulatero2004increased}.

\section*{Why This Test Is Ordered}
\begin{itemize}
  \item Screening for primary aldosteronism in hypertensive patients with: resistant hypertension (BP uncontrolled on 3 medications including a diuretic); hypertension with spontaneous or diuretic-induced hypokalemia; hypertension with adrenal incidentaloma; early-onset hypertension (<30 years); severe hypertension (systolic >160 mmHg or diastolic >100 mmHg); or hypertension with a family history of early hypertension or cerebrovascular accident before age 40
  \item Diagnostic workup of hypertension with hypokalemia (<3.5 mmol/L), especially if spontaneous (not diuretic-induced) and accompanied by metabolic alkalosis
  \item Investigation of adrenal incidentaloma >1 cm to assess for aldosterone production
  \item Diagnosis of primary aldosteronism using the aldosterone-to-renin ratio (ARR) as a screening test, followed by confirmatory testing (saline suppression, captopril challenge, oral sodium loading, or fludrocortisone suppression)
  \item Lateralization of aldosterone excess in confirmed primary aldosteronism using adrenal vein sampling (AVS) with aldosterone and cortisol measurement
  \item Monitoring response to mineralocorticoid receptor antagonist therapy (spironolactone, eplerenone) or post-adrenalectomy in primary aldosteronism
  \item Diagnosis of congenital adrenal hyperplasia due to 11-beta-hydroxylase or 17-alpha-hydroxylase deficiency (low aldosterone with low cortisol, elevated precursors)
  \item Investigating hypoaldosteronism (Type IV renal tubular acidosis) in patients with hyperkalemia, metabolic acidosis, and mild renal insufficiency typical of diabetic nephropathy
\end{itemize}

\section*{Is This a Primary or Secondary Test}
The aldosterone-to-renin ratio (ARR) is a primary screening test for primary aldosteronism in at-risk hypertensive populations. However, the ARR is not a general screening test for all hypertensive patients because the prevalence of primary aldosteronism (5--10\%) does not justify universal screening. The Endocrine Society (2016) recommends screening in high-prevalence groups \cite{funder2016management}. Isolated aldosterone measurements without corresponding renin have limited clinical utility because aldosterone must be interpreted in the context of renin activity (volume status and RAAS activation). Confirmatory tests (saline suppression, captopril challenge) are required after a positive ARR before the diagnosis of primary aldosteronism is confirmed \cite{funder2016management}.

\section*{How This Test Is Performed}
A healthcare professional draws a blood sample from a vein into an EDTA (purple-top) tube or a serum separator tube (gold-top), depending on whether renin is being measured simultaneously (EDTA preferred for renin activity). The draw should occur mid-morning (8--10 AM) after the patient has been upright (sitting, standing, walking) for at least 2 hours, with the patient seated for 5--15 minutes before venipuncture. Aldosterone is relatively stable at room temperature for 24 hours and at 2--8 degrees C for 7 days. For the saline suppression test (confirmatory): the patient lies supine for at least 1 hour, then receives 2 liters of 0.9\% normal saline IV over 4 hours. Aldosterone, renin, cortisol, and potassium are measured at baseline and after infusion. Failure to suppress aldosterone <5--10 ng/dL (assay-dependent) confirms primary aldosteronism. For adrenal vein sampling (AVS): catheters are placed in both adrenal veins via femoral vein access by interventional radiology. Aldosterone and cortisol are measured simultaneously from both adrenal veins and a peripheral vein (often inferior vena cava), both before and after cosyntropin (ACTH 1-24) stimulation. A cortisol-corrected aldosterone ratio (aldosterone/cortisol) from the dominant side >4 times the non-dominant side confirms unilateral disease amenable to surgery.

\section*{What Preparation Is Needed}
Fasting for 8--10 hours before testing is recommended but not mandatory. The patient should be on a liberal sodium diet for at least 3 days before testing (sodium restriction stimulates RAAS and can cause false-positive ARR). Hypokalemia must be corrected before testing because low potassium suppresses aldosterone secretion, producing false-negative results (target potassium >4.0 mmol/L with oral KCl supplementation). Several antihypertensive medications must be discontinued or substituted before ARR testing: mineralocorticoid receptor antagonists (spironolactone, eplerenone) --- hold for at least 4--6 weeks; high-dose amiloride and triamterene --- hold for 2--4 weeks; ACE inhibitors, ARBs, and direct renin inhibitors (aliskiren) --- these can elevate renin and lower ARR, causing false-negatives; they should be held for 2 weeks if possible and substituted with a non-interfering agent (verapamil slow-release, hydralazine, prazosin, doxazosin); beta-blockers, clonidine, and alpha-methyldopa --- suppress renin, causing false-positive ARR, and should be held for 2 weeks if possible; diuretics --- stimulate renin and should be held for 4 weeks. The patient should remain upright (sitting, standing, walking) for at least 2 hours before the blood draw, then sit for 5--15 minutes before venipuncture, because upright posture stimulates renin and aldosterone. NSAIDs suppress renin and should be held for 2 weeks if possible. Oral contraceptives and estrogen therapy may decrease direct renin concentration (not renin activity) and should be documented.

\section*{Contraindications and Precautions}
Factors falsely lowering the ARR (false-negatives): hypokalemia (<3.5 mmol/L) suppresses aldosterone --- potassium must be corrected before testing; sodium restriction and diuretic use stimulate renin disproportionately, lowering ARR; pregnancy (renin and aldosterone physiologically elevated by progesterone, decreasing ARR); ACE inhibitors, ARBs, and direct renin inhibitors increase renin, decreasing ARR; renovascular hypertension markedly elevates renin, decreasing ARR; malignant hypertension. Factors falsely elevating the ARR (false-positives): beta-blockers, clonidine, alpha-methyldopa suppress renin more than aldosterone; advanced age (>65 years) --- renin declines with age, and ARR increases; chronic kidney disease (decreased renin in CKD stages 3--5); African ancestry (tendency toward low-renin hypertension); high-sodium diet suppresses renin; potassium loading. The ARR can be affected by the time of day, posture, sodium balance, and potassium balance. The reproducibility of the ARR is modest (30--50\% of patients with initially positive ARR will have a negative ARR on repeat testing), so a single abnormal ARR should be confirmed by repeat testing before proceeding to confirmatory testing.

\section*{Risks}
From venipuncture: slight pain, bruising, hematoma. From saline suppression test: volume expansion in a hypertensive patient with increased aldosterone may worsen hypertension transiently; patients with heart failure or renal insufficiency may not tolerate 2 L normal saline and this test is contraindicated in these groups. From adrenal vein sampling: groin hematoma (common), contrast reaction, adrenal vein rupture (rare, 0.2--1\%, can cause adrenal hemorrhage), adrenal vein thrombosis (rare), and radiation exposure. The right adrenal vein is technically challenging to cannulate (success rate 70--95\% depending on operator experience). From captopril challenge test: hypotension (captopril 25--50 mg orally reduces BP). Contraindicated in pregnancy.

\section*{What Happens After the Test}
Aldosterone results are available within 1--4 hours on automated analyzers. Interpretation of the ARR: a positive screen (ARR >20--30 with PRA, or >2--3.7 with DRC) combined with aldosterone >10--15 ng/dL prompts confirmatory testing. If confirmatory testing is positive (failure to suppress aldosterone), primary aldosteronism is confirmed. The next step is CT adrenal imaging (non-contrast fine cut, 1--3 mm slices) to evaluate for adrenal adenoma versus bilateral hyperplasia and to exclude adrenal carcinoma (>4--6 cm, heterogeneous). Adrenal vein sampling is required in patients who are surgical candidates to determine whether aldosterone excess is unilateral (adenoma, amenable to laparoscopic adrenalectomy) or bilateral (hyperplasia, managed medically with mineralocorticoid receptor antagonists) \cite{rossi2014expert}. Surgery is contraindicated without AVS confirmation of lateralization because CT may show a non-functioning incidentaloma while the true aldosterone source is the contralateral side (discordance rate 20--40\%) \cite{rossi2014expert}.

\section*{Understanding Results}

\subsection*{Below Normal}
\begin{itemize}
  \item \textbf{Low aldosterone with elevated renin (low ARR):} Primary adrenal insufficiency (Addison disease). Destruction of all three adrenal cortical zones including zona glomerulosa impairs aldosterone production. Hyponatremia, hyperkalemia, metabolic acidosis, and volume depletion occur. Requires fludrocortisone replacement (0.05--0.2 mg/day) and hydrocortisone.
  \item \textbf{Low aldosterone with elevated renin and hyperkalemia:} Hypoaldosteronism. The most common form is Type IV renal tubular acidosis (hyporeninemic hypoaldosteronism), typically in older adults with mild-to-moderate diabetic nephropathy or chronic interstitial nephritis. Renin is low because of juxtaglomerular apparatus damage, leading to decreased aldosterone despite hyperkalemia stimulation.
  \item \textbf{Low aldosterone with low or normal renin:} 11-beta-hydroxylase or 17-alpha-hydroxylase deficiency (rare CAH variants) --- impaired cortisol and aldosterone synthesis with accumulation of DOC (deoxycorticosterone), which acts as a mineralocorticoid causing hypertension and hypokalemia (low aldosterone levels but hypertension paradoxically). Also seen in Liddle syndrome (constitutively active ENaC, suppressed aldosterone and renin with hypertension, hypokalemia, metabolic alkalosis), and apparent mineralocorticoid excess (11-beta-HSD2 deficiency preventing cortisol-to-cortisone conversion, cortisol activating MR).
  \item \textbf{Suppressed aldosterone on confirmatory testing:} Aldosterone suppressed to <5--10 ng/dL (assay-dependent) after saline infusion or to <8 ng/dL after captopril --- excludes primary aldosteronism. Essential hypertension is the most common cause of low-renin, normal-aldosterone hypertension with a normal ARR.
\end{itemize}

\subsection*{Above Normal}
\begin{itemize}
  \item \textbf{Primary aldosteronism (Conn syndrome):} Elevated aldosterone (>15 ng/dL, often >20 ng/dL) with suppressed renin (PRA <1 ng/mL/hr or DRC <10 mU/L) and elevated ARR (>30 with PRA, >3.7 with DRC). Spontaneous hypokalemia in only 30--40\% of patients (most are normokalemic). Confirmatory testing positive (failure to suppress aldosterone after saline or captopril). Aldosterone-producing adenoma (APA) accounts for 30--40\% --- unilateral, >1 cm, treated by laparoscopic adrenalectomy, which cures hypertension in 30--60\% and improves it in most. Bilateral idiopathic hyperaldosteronism (IHA) accounts for 60--70\% --- bilateral micronodular hyperplasia, treated medically with spironolactone or eplerenone, surgery not indicated.
  \item \textbf{Secondary hyperaldosteronism (elevated aldosterone with elevated renin, normal ARR):} Renovascular hypertension (atherosclerotic renal artery stenosis or fibromuscular dysplasia) --- renal hypoperfusion drives excessive renin and aldosterone secretion. Renin-secreting juxtaglomerular cell tumors (reninomas). Edematous states: congestive heart failure, cirrhosis with ascites, nephrotic syndrome --- effective circulating volume depletion activates RAAS. Diuretic therapy (thiazides, loop diuretics). Bartter syndrome and Gitelman syndrome (renal salt-wasting with secondary hyperaldosteronism, hypokalemic metabolic alkalosis, normal or low blood pressure). Malignant hypertension (accelerated RAAS activation).
  \item \textbf{Familial hyperaldosteronism Type I (glucocorticoid-remediable aldosteronism, GRA):} Caused by a chimeric CYP11B1/CYP11B2 gene that places aldosterone synthase under ACTH control rather than RAAS control. Autosomal dominant, onset in childhood/young adulthood, family history of early-onset hypertension and hemorrhagic stroke. Aldosterone suppressed by dexamethasone (diagnostic). Genetic testing confirms the chimeric gene.
  \item \textbf{Familial hyperaldosteronism Type III:} Germline mutations in the KCNJ5 potassium channel gene causing constitutive aldosterone production. Severe hypertension, massive bilateral adrenal hyperplasia, often requires bilateral adrenalectomy. Rare.
  \item \textbf{Artifactual ARR elevation:} Beta-blocker therapy (suppresses renin), advanced age (physiological renin decline), chronic kidney disease, low-sodium diet, pregnancy (third trimester ARR may be elevated physiologically).
\end{itemize}

\section*{Normal Results}
\begin{itemize}
  \item \textbf{Serum/plasma aldosterone (adults, supine):} 1--16 ng/dL (30--440 pmol/L; Roche, Siemens); upright: 4--31 ng/dL (110--860 pmol/L) --- upright aldosterone is 2--4 times higher than supine
  \item \textbf{Plasma renin activity (PRA, supine):} 0.2--1.6 ng/mL/hr; upright (2 hr ambulation): 1.5--5.2 ng/mL/hr
  \item \textbf{Direct renin concentration (DRC, supine):} 3--16 mU/L (males), 2--16 mU/L (females); upright: 10--40 mU/L
  \item \textbf{Aldosterone-to-renin ratio (ARR) using PRA:} <20 ng/dL per ng/mL/hr (negative screen); >30 suggests primary aldosteronism
  \item \textbf{Aldosterone-to-renin ratio (ARR) using DRC:} <2.0 (pmol/L per mU/L, negative screen); >3.7 suggests primary aldosteronism
  \item \textbf{24-hour urinary aldosterone (normal sodium diet):} 2--20 mcg/24 hours (5.5--55 nmol/24 hours)
  \item \textbf{Saline suppression test:} Aldosterone <5--10 ng/dL post-2L NS (excludes primary aldosteronism)
  \item \textbf{Captopril challenge:} Aldosterone suppressed >30\% from baseline to <8 ng/dL (excludes primary aldosteronism)
\end{itemize}

\section*{Follow-up Tests}
\begin{itemize}
  \item \textbf{Plasma renin activity (PRA) or direct renin concentration (DRC):} Always measured simultaneously with aldosterone. ARR interpretation is impossible without renin. PRA measures enzymatic activity; DRC measures immunoreactive renin mass. DRC is simpler (automated) but less validated than PRA in the literature.
  \item \textbf{Saline infusion test:} Confirmatory test. IV 0.9\% saline 2 L over 4 hours, aldosterone measured at baseline and after infusion. Aldosterone <5--10 ng/dL (depending on assay) excludes primary aldosteronism. Sensitivity 88\%, specificity 85--90\%.
  \item \textbf{Captopril challenge test:} Captopril 25--50 mg orally; aldosterone, renin, cortisol measured at baseline and 2 hours. Normal: aldosterone suppressed >30\%. Primary aldosteronism: aldosterone remains elevated >12--15 ng/dL with suppressed renin, ARR remains elevated.
  \item \textbf{Oral sodium loading test:} High-sodium diet (5000 mg/day for 3 days) plus fludrocortisone 0.1 mg every 6 hours for 4 days or NaCl tablets 2 g three times daily. 24-hour urine aldosterone >12--14 mcg on high-sodium diet confirms primary aldosteronism.
  \item \textbf{Fludrocortisone suppression test:} Fludrocortisone 0.1 mg every 6 hours for 4 days with high-sodium diet. Aldosterone >6 ng/dL at 10 AM on day 4 confirms primary aldosteronism.
  \item \textbf{Adrenal CT (fine-cut, 1--3 mm):} Detects adrenal adenoma (>1 cm, low attenuation <10 HU on non-contrast CT, rapid washout >50\% on 15-minute delayed images). Adrenal carcinoma >4--6 cm, irregular, heterogeneous, high attenuation. Cannot determine functionality of adenoma.
  \item \textbf{Adrenal vein sampling (AVS):} Gold standard for lateralization. Aldosterone/cortisol ratio from each adrenal vein compared. Lateralized ratio >4:1 (dominant:non-dominant) with contralateral suppression suggests unilateral adenoma treatable by surgery. Performed with cosyntropin stimulation to increase gradient. Discordance between CT and AVS occurs in 20--40\%, making AVS essential for surgical candidates.
  \item \textbf{Genetic testing for familial hyperaldosteronism:} CYP11B1/CYP11B2 chimeric gene for GRA (Type I). KCNJ5, CACNA1D, CACNA1H, CLCN2 gene sequencing for other familial forms.
  \item \textbf{24-hour urinary aldosterone and potassium:} Elevated urinary aldosterone with inappropriately high urinary potassium (renal potassium wasting) supports primary aldosteronism.
\end{itemize}

\section*{References}
\bibliographystyle{unsrtnat}
\bibliography{references}

\clearpage
